\section*{Abstract}
This project builds and tests a system of three models for reviewing contracts automatically, using CUAD (510 contracts, 41 clause categories) and a held-out test set made of 20\% of the contracts. The system finds which clause types a contract contains, extracts the clause text, gives each clause type a High, Medium or Low risk level with a plain-English line, and runs as a working app. In our first setup, a DistilBERT model that reads the contract in windows lost to a simple TF--IDF model that reads the whole document (micro-F1 0.692 against 0.775), and the gap survived changes to the training data, seed, pooling rule and threshold. We then found that the model read only about half of each window. Once it reads the whole window and both models are tuned on the same validation split, the transformer ties the tuned baseline (0.797 against 0.785), an ensemble of the two beats it (0.809, with a 95\% interval for the lead of $[+0.011, +0.037]$), and a recall-first setting finds 90.3\% of High-risk clauses, against 63.1\% for the first setup's best ensemble. Extracting the clause text reaches a token-F1 of 0.764 when the model is given the right window. The fine-tuned summarizer outputs one of 41 sentences we wrote ourselves, word for word, for every input, and scores 0.116 against clause-specific references. A calibration map fitted on the held-out contracts lowers the confidence score's calibration error from 0.206 to 0.016. Over the whole system only 22.1\% of real clauses get through all three steps: the presence model now finds the right window 82.6\% of the time, but the span model then succeeds in only 28.0\% of those cases.

\vspace{1cm}
\textbf{Keywords:} legal NLP; CUAD; clause extraction; category-level risk prioritization; contract summarization; held-out evaluation
\pagebreak
